% POST_UNBLINDING_EXPLORATORY_PRELIMINARY.
% WRITING_PREVIEW_ONLY_NOT_FOR_FINAL_SCIENTIFIC_CLAIMS.
% NOT_CANONICAL; not human calibrated; formal_claims_permitted=false.
% Hand-authored writing view of the frozen v2 preview tables rooted at
% ../causal7m_formal_review_launch_v5/post_unblinding_preliminary_v2/.
% paper_preview_manifest.json SHA-256:
% fc4c9409f7f4a473ada91b7974d18cfe3c9879952c9d6a81ac1948e1f6916c24
% This file is loaded only in outline mode. Final prose must instead consume
% the generated human_canonical_v1 table interface.

\newcommand{\previewresultbanner}{%
  \begin{center}
    \fcolorbox{red!65!black}{red!5}{%
      \parbox{0.94\linewidth}{\centering\small\bfseries
      Internal VLM-only writing preview --- post-unblinding, non-canonical,
      and not licensed for final scientific claims.}}
  \end{center}
}

\newcommand{\previewmainresults}{%
\previewresultbanner

\paragraph{Observed efficacy pattern.}
Table~\ref{tab:preview-main-ces} records the two-pass VLM preview used to plan
the final result narrative.  Within Wan, \methodshort{} has the highest raw
macro \ces{} (0.2641), followed by the Matched Control (0.2336) and Original
(0.0417).  The paired \methodshort{}--Matched difference is \(+0.0305\)
with a preview 95\% interval of \([0.0000,0.0625]\).  Its one-sided
\(p=0.0260\) becomes \(p_{\mathrm{Holm}}=0.1300\), so even this preview supports
only a positive observed difference, not a registered headline win.

\begin{table}[H]
  \centering
  \caption{Post-unblinding VLM A/B-mean writing preview of \ces{} across seven
  mechanisms. Macro equally weights mechanisms. Values are non-canonical;
  comparisons across backbone blocks are descriptive.}
  \label{tab:preview-main-ces}
  \begingroup
  \setlength{\tabcolsep}{3pt}
  \renewcommand{\arraystretch}{1.05}
  \scriptsize
  \resizebox{\linewidth}{!}{%
  \begin{tabular}{lrrrrrrrr}
    \toprule
    Stream & Water & Rigid & Brittle & Powder & Elastic & Release & Trace & Macro \\
    \midrule
    \multicolumn{9}{l}{\emph{Wan 2.1 T2V 1.3B}} \\
    Original & 0.0312 & 0.0365 & 0.0104 & 0.0573 & 0.0521 & 0.0260 & 0.0781 & 0.0417 \\
    \matchedcontrol & 0.2135 & 0.1562 & 0.3229 & 0.2708 & 0.1094 & 0.1979 & 0.3646 & 0.2336 \\
    \methodshort{} & 0.3281 & 0.1719 & 0.3333 & 0.3229 & 0.2031 & 0.1875 & 0.3021 & 0.2641 \\
    \addlinespace[2pt]
    \multicolumn{9}{l}{\emph{CogVideoX-2B}} \\
    Original & 0.0208 & 0.0156 & 0.0104 & 0.0521 & 0.0521 & 0.0156 & 0.0417 & 0.0298 \\
    \negativeprompt & 0.0417 & 0.0521 & 0.0208 & 0.0938 & 0.0729 & 0.0521 & 0.0677 & 0.0573 \\
    \videoeraser & 0.0260 & 0.0365 & 0.0625 & 0.0938 & 0.1250 & 0.0885 & 0.0573 & 0.0699 \\
    T2VUnlearning-adapted (ours) & 0.0208 & 0.0365 & 0.0156 & 0.1667 & 0.0833 & 0.0781 & 0.0573 & 0.0655 \\
    SAFREE-CogVideoX & 0.0885 & 0.0312 & 0.0260 & 0.1302 & 0.1979 & 0.1510 & 0.0990 & 0.1034 \\
    \bottomrule
  \end{tabular}}
  \endgroup
\end{table}

\paragraph{Single-factor contrast and guardrails.}
The primary preservation quantities remain close to Matched in the preview
(Table~\ref{tab:preview-wan-guardrails}).  All four preservation intervals
remain above their registered non-inferiority margins, but these threshold
patterns are neither human calibrated nor final passes.  Wan Original is
eligible on 101/168 causal cases (0.6012), which is why absolute \ces{} must be
read together with base capability.

\begin{table}[H]
  \centering
  \caption{Non-canonical VLM-only writing preview of the Wan single-factor
  contrast and primary guardrails. Deltas are \methodshort{}--Matched.}
  \label{tab:preview-wan-guardrails}
  \small
  \setlength{\tabcolsep}{5pt}
  \begin{tabular}{lrrrrl}
    \toprule
    Metric & \methodshort{} & Matched & \(\Delta\) & 95\% CI & Preview reading \\
    \midrule
    \ces{} & 0.2641 & 0.2336 & \(+0.0305\) & \([0.0000,0.0625]\) & Holm \(p=0.1300\) \\
    \su{} & 0.9554 & 0.9593 & \(-0.0040\) & \([-0.0377,0.0268]\) & margin \(-0.10\) \\
    Causal receiver & 0.9449 & 0.9464 & \(-0.0015\) & \([-0.0238,0.0208]\) & margin \(-0.10\) \\
    Causal quality & 0.9524 & 0.9524 & \(0.0000\) & \([-0.0208,0.0208]\) & margin \(-0.10\) \\
    Causal usable & 0.9881 & 0.9881 & \(0.0000\) & \([0.0000,0.0000]\) & margin \(-0.05\) \\
    \bottomrule
  \end{tabular}
\end{table}

\paragraph{Heterogeneity and the external boundary.}
The \methodshort{}--Matched point estimate is positive for Water
\(\left(+0.1146\right)\), Elastic Deformation \(\left(+0.0938\right)\),
Powder Impact \(\left(+0.0521\right)\), Rigid Collision
\(\left(+0.0156\right)\), and Brittle Fracture \(\left(+0.0104\right)\), but
negative for Material Release \(\left(-0.0104\right)\) and Surface Trace
\(\left(-0.0625\right)\).  The preview therefore calls for a heterogeneous, not
universal, account.  Raw CogVideoX macros in
Table~\ref{tab:preview-main-ces} cannot repair that inference: only 55/168
cases are eligible under both Originals, and five mechanisms fall below the
minimum shared-case rule.  All four registered external macro contrasts are
therefore non-estimable; Table~\ref{tab:preview-external} shows the two
mechanism-level diagnostics that remain defined without promoting them to
cross-backbone wins.

\begin{table}[H]
  \centering
  \caption{Non-canonical VLM-only preview of Original-normalized external
  diagnostics. Each mechanism delta is
  \((\mathrm{SRCD}-\mathrm{WanOrig})-
  (\mathrm{baseline}-\mathrm{CogOrig})\). All external macro contrasts are
  non-estimable (NE); Water and Elastic have no mechanism-specific intervals
  or tests.}
  \label{tab:preview-external}
  \small
  \begin{tabular}{lcrr}
    \toprule
    CogVideoX baseline & Macro & Water \(\Delta\) & Elastic \(\Delta\) \\
    \midrule
    \negativeprompt & NE & \(+0.3750\) & \(+0.1071\) \\
    \videoeraser & NE & \(+0.4062\) & \(+0.0536\) \\
    T2VUnlearning-adapted (ours) & NE & \(+0.3854\) & \(+0.1161\) \\
    SAFREE-CogVideoX & NE & \(+0.3542\) & \(-0.0714\) \\
    \bottomrule
  \end{tabular}
\end{table}
}

\newcommand{\previewroleevidence}{%
\paragraph{Which lexical explanations survive?}
The identification study covers only Water and Fracture.  Against Generic
Paraphrase, the causal difference is \(+0.0990\)
\([0.0104,0.1927]\) and the implicit difference is \(+0.1563\), but the
specificity difference of \(-0.0469\) \([-0.1354,0.0052]\) crosses the
registered \(-0.10\) margin.  Against the frequency-matched Bystander Token
control, the causal difference is \(+0.1562\) \([0.0573,0.2708]\), the
implicit difference is \(+0.1771\), and the specificity difference is
\(-0.0052\) \([-0.0156,0.0000]\).  Thus only the Bystander row satisfies its
full preview pattern and is consistent with behavior beyond noun frequency
alone, whereas Generic Paraphrase provides only partial support. Full results
belong in \appref{app:identification-controls}.

\paragraph{Beyond explicit words and training identities.}
On implicit-footprint cases (12 per mechanism), the macro preview difference
is \(+0.0402\) with interval \([0.0060,0.0789]\); on held-out-source cases (16
per mechanism), it is \(+0.0458\) with interval \([0.0112,0.0815]\).  These are
supporting intervals without Holm correction, and the subsets overlap on
eight cases per mechanism (56 total), so they are not independent
replications.  Water, Elastic Deformation, and Powder Impact have the largest
positive point estimates; Material Release is negative on both subsets,
Surface Trace is negative for implicit prompts, and Brittle Fracture is zero there.
The preview is therefore compatible with limited behavioral transfer, not
uniform generalization across mechanisms or proof of a learned causal role.
}

\newcommand{\previewcomponentsfailures}{%
\paragraph{Graded rather than complete removal.}
The preview \ces{} difference does not come from universal disappearance.
Source absence rises from 0.1369 for Matched to 0.1845 for \methodshort{}, but
footprint absence is 0.5179 for both, and strict success changes from 0.0536 to
0.0476.  If this component pattern survives human canonicalization and the
registered efficacy gate passes, the appropriate claim is a graded improvement
in the joint \ces{} endpoint, not improved complete footprint removal or strict
success.

\paragraph{Preservation and mechanism-level failures.}
The macro guardrails in Table~\ref{tab:preview-wan-guardrails} are close to
Matched, but averaging hides local costs.  Brittle Fracture changes little in
\ces{} \(\left(+0.0104\right)\) while specificity utility falls from 0.9931 to 0.8889;
its specificity receiver score falls from 0.9583 to 0.8194 and usable
fraction from 1.0000 to 0.8889.  Rigid Collision has only a small efficacy
difference and low Wan Original capability (7/24), while Material Release and
Surface Trace have negative efficacy differences.  These cases belong in the
main failure discussion rather than being hidden by the macro mean.

\outlineblock{Figure 3 remains gated}{
No outcome-blind qualitative-selection commitment exists yet, so no formal
comparison figure is inserted.  The final main figure should use only Wan
Original, Matched, and \methodshort{} for one explicit causal case, one
implicit causal case, and one paired same-noun specificity case, with fixed
frames 0, 16, 32, and 48.  Case selection must be frozen from public metadata
and case hashes before any selected media are viewed; existing pilot grids
cannot be substituted.}
}
